% !TEX root = userguide.tex

\chapter{Some examples using Galerkin \fem: the Poisson and Stokes equations}
\pgst{tfemch}

\label{ch:examples}

\def\chaptertitle{Some examples}
\def\headdate{4th December 2025}

In this chapter some examples for the Poisson and Stokes equations using
Galerkin \fem\ will be given.
The first example will be fully
explained line by line in order to get a good introduction to \tfem\
and `learn by example'. The succeeding examples
will be explained more briefly.

\section{The Poisson equation}

In this section some examples using the Poisson equation will be given, but
first some theory.

\subsection{Weak form of the Poisson equation. Galerkin \fem}
\label{sec:poissonweak}

We consider the Poisson equation in a region $\Omega$:
\begin{equation}
     -\nabla^2 u = f\text{ on }\Omega
   \label{eq:poisson}
\end{equation}
with boundary conditions
\begin{gather}
    u = u_D\text{ on }\Gamma_D\\
    -\pderiv{u}{n}=-\vec n\cdot\nabla u=h_{\text{N}}\text{ on
}\Gamma_{\text{N}}
\end{gather}
where the boundary $\Gamma=\Gamma_D\cup\Gamma_{\text{N}}$ has been split into a
Dirichlet and a Neumann part.

The weak form can be obtained by multiplying Eq.~\eqref{eq:poisson}
with a test function $v$:
\[
   (v,-\nabla^2 u-f)=0\qquad\text{for all }v
\]
where the standard inner product on $L^2(\Omega)$ is
\[
    (a,b)=\int_\Omega ab \D x
   \label{eq:innerproduct}
\]
With
\[
   \nabla\cdot\big(\nabla u\big)v=\nabla\cdot\big((\nabla u)v)\big)-\nabla
u\cdot\nabla v
\]
and the divergence theorem (Gauss):
\[
    \int_\Omega\nabla\cdot\vek a\D x=\int_\Gamma\vek n\cdot\vek a\D s
\]
we get the weak form: Find $u\in S$ such that
\begin{equation}
   (\nabla v,\nabla u) +(v,h_{\text{N}})_{\Gamma_{\text{N}}}=(v,f)
   \label{eq:weakformpoisson}
\end{equation}
for all $v\in V$, where $S$ and $V$ are suitable functional spaces for $u$ and
$v$, respectively. Notes:
\begin{itemize}
  \itemsep -1ex
   \item $v=0$ on $\Gamma_D$, $u=u_D$ on $\Gamma_D$
   \item we have silently introduced:
\begin{align*}
   (\vek a, \vek b)&=\int_\Omega\vek a\cdot\vek b\D x\\
   (a,b)_{\Gamma}&=\int_\Gamma a b\D s
\end{align*}
\end{itemize}

The weak form can be used to obtain an approximate solution using the Galerkin \fem\
technique. For this we define approximating 
spaces $S_h$ and $V_h$, or in other words, we divide the region into elements:
\[
\Omega=\cup_e\Omega_e, \quad\Omega_{e}\cap\Omega_{f}=\emptyset\quad
   \text{for }e\neq f
\]
and interpolate $u$ and $v$ by
polynomials on each element (approximation denoted by $u_h$ and $v_h$):
\begin{align*}
    u_h(\vek x,t)&=\tsum_i u_i(t)\phi_i(\vek x)=\col\phi^T(\vek x)\col u(t)
\\
    v_h(\vek x,t)&=\tsum_i v_i(t)\phi_i(\vek x)=\col\phi^T(\vek x)\col v(t)
\end{align*}
where $\col\phi(\vek x)$ are global shape functions.  Substituting into the
weak form leads to
\[
 \col v^T\mat K\col u=
        \col v^T\col f \quad\text{for all }\col v
\]
or
\[
 \mat K\col u= \col f
\]
where the `stiffness matrix' $\mat K$ and
the right-hand side $\col f$ are given by
\begin{alignat*}{2}
 \mat K&=(\nabla{\col\phi},\nabla{\col\phi^T})\qquad\qquad\qquad\text{  or}&\qquad\qquad
K_{ij}&=(\nabla{\phi_i}, \nabla{\phi_j}) \\[2ex]
\col f&= (\col\phi,f)-(\col\phi,h_{\text{N}})_{\Gamma_{\text{N}}} &
f_i&= (\phi_i,f)-(\phi_i,h_{\text{N}})_{\Gamma_{\text{N}}}
\end{alignat*}
where the left column is in matrix notation and the right column in index
notation.



\subsection{A Poisson problem on a unit square with Dirichlet boundary 
            conditions}
\label{sec:poisson1}

We consider a Poisson problem:
\begin{equation}
     -\nabla^2 u = f
\end{equation}
on a unit square $(x,y)\in(0,1)\times(0,1)$ (see figure~\ref{fig:unitsquare}).
We substitute the following solution for $u$:
\begin{equation}
     u=1+\cos(\pi x)\cos(\pi y)
     \label{eq:exactu}
\end{equation}
and find the corresponding $f$:
\begin{equation}
     f=2\pi^2\cos(\pi x)\cos(\pi y)
     \label{eq:exactf}
\end{equation}
\begin{figure}[ht!]
   \begin{center}
   \includegraphics{unitsquare}
   \end{center}
   \caption{Unit square and the points and curves defined by the meshgenerator
      \texttt{quadrilateral2d}. Note the default directions of the curves.}
    \label{fig:unitsquare}
\end{figure}

 In the following we assume
that, as suggested in Sec~\ref{sec:gettingstarted}, the user performed the command 
\begin{quote}
   \texttt{getexample poisson}
\end{quote}
and all the necessary steps before the
\texttt{make}-command step.

The functions given by Eq.~\eqref{eq:exactu} and Eq.~\eqref{eq:exactf}
need to be programmed in a module. The first part of the file
\texttt{poisson\_functions.f90} might look something like:

\begin{listing}{1}
module poisson_functions_m
                                                                                
  use kind_defs_m
                                                                                
  implicit none
                                                                                
contains

  function func ( nr, x )
    integer, intent(in) :: nr
    real(dp), intent(in), dimension(:) :: x
    real(dp) :: func

    real(dp) :: pi

    pi = 4 * atan(1._dp)

    select case(nr)
      case(1)
        func = 0
      case(2)
        func = 1
      case(3)
        func = 1 + cos(pi*x(1))*cos(pi*x(2))
      case(4)
        func = 2*pi**2*cos(pi*x(1))*cos(pi*x(2))
      case(5)
        func = -pi*sin(pi*x(1))*cos(pi*x(2))
      case(6)
        func = -pi*cos(pi*x(1))*sin(pi*x(2))
      case default
        write(*,'[a,i0]') 'Error func: wrong function number: ', nr
        stop
    end select

  end function func

  ...

end module poisson_functions_m

\end{listing}

Some explanations:
\begin{description}
   \item [line 1] In \tfem, modules must be in separate files or be stuffed in
front of the main program (not recommended). Furthermore, the filename must be
\texttt{<basename>.f90}, where \texttt{<basename>} is either the module name or
the module name without the \texttt{\_m}. These are requirements for the script
\texttt{fmm} to work properly\footnote{In \texttt{fmm} it is allowed to use
\texttt{<basename>.f} (fixed format source), but this should be limited to
source code which is copy/pasted mostly from old style Fortran.}. Throughout
\tfem\ the convention of extending a module name with \texttt{\_m} is used.

   \item [line 3] The module \texttt{kind\_defs\_m} defines the kind parameters
    for precision, among which \texttt{dp}: the precision that used to be
    `double precision' in old-style Fortran. 

    \item [lines 23--26] For \texttt{nr=3} we have $u$ and for \texttt{nr=4} we have
    $f$.
\end{description}

Now we have to write a main program in \ftn\ to solve the problem. A global
overview of a typical main program is depicted in Fig.~\ref{fig:mainprogram}.
The main program doesn't have to be sequential and all
the programming facilities of \ftn\
are available. For example, in Fig.~\ref{fig:mainprogram} two possible loops
are illustrated: (1) time-stepping scheme, (2) an adaptive mesh scheme.
\begin{figure}[ht!]
   \begin{center}
   \includegraphics{main}
   \end{center}
   \caption{Schematic overview of a main program}
    \label{fig:mainprogram}
\end{figure}


The main program of the first Poisson
problem is in the file \texttt{poisson1.f90}.

The Poisson element routines
in the add-on `convection diffusion' (see Sec.~\ref{sec:convdiff_elements}) are used for
the building of the system. 

We will give an explanation of the main program line by line.

\begin{listing}{1}
! Poisson problem on a unit square with Dirichlet boundary conditions

program poisson1

  use kind_defs_m
  use mesh_m
  use meshgen_m
  use problem_m
  use system_m
  use hsl_ma57_m
  use poisson_elements_m
  use poisson_functions_m
  use figplot_m

\end{listing}
Here all modules needed are included using the \texttt{use} statement:
\begin{description}
   \item[line 6] this defines the mesh type (\texttt{mesh\_t}) and more.
   \item[line 7] this is the mesh generation module, which includes a simple
                 mesh generator and various other routines for mesh generation.
   \item[line 8] this defines the problem type (\texttt{problem\_t}) and more.
   \item[line 9] this defines the type of the
                 vectors and matrices for the system and many
                 routines, such as for building the system.
   \item[line 10] this defines the types and associated routines to run the HSL
                  sparse direct solver MA57.
   \item[line 11] this defines the element routines for the Poisson equation.
                  These routines are defined in a module, which can be found in
                  the source file 
                  \texttt{poisson\_elements.f90}. It is
                  not part of the core of \tfem but part of the add-on 
                  `convection diffusion' (see Sec.~\ref{sec:convdiff_elements},
                   which must
                  have already been compiled into to the \tfem-library.
   \item[line 12] this includes the user function \texttt{func} for use in the main program.
   \item[line 13] this includes some plot routines.
                  This module is part of an add-on
                  (see Sec.~\ref{sec:figplot}), which must
                  have already been compiled into to the \tfem-library.
\end{description}
The first five use statements appear in almost every \tfem-program. Therefore
these modules (and other standard modules) have been included in a module
\texttt{tfem\_m}. Therefore the statements above could have been replaced by
\begin{verbatim}
  program poisson1

  use tfem_m
  use hsl_ma57_m
  use poisson_elements_m
  use poisson_functions_m
  use figplot_m
 
\end{verbatim}
which is what we did in all subsequent examples.

\begin{listing}{14}
  implicit none


! constants

  integer, parameter :: &
    uintpl = 8,         & ! scalar interpolation
    gauss = 3,          & ! 3x3 integration of quads
    gaussb = 3,         & ! 3 point integration of boundary elements
    nx=20,              & ! number of elements in x
    ny=20,              & ! number of elements in y
    funcnr=4              ! function number for the right-hand side

  real(dp), parameter :: &
    alpha = 1._dp     ! diffusion coefficient


! definitions

  type(meshgen_options_t) :: meshgen_options
  type(mesh_t) :: mesh
  type(input_probdef_t) :: input_probdef
  type(problem_t) :: problem
  type(sysmatrix_t) :: sysmatrix
  type(sysvector_t) :: sol, rhsd, solexact
  type(plot_options_t) :: plot_options
  type(coefficients_t) :: coefficients

\end{listing}
Here all the variables are defined that are part of the main program:
\begin{description}
   \item[line 14] every \ftn\ program should contain this statement, which
       ensures that all variables are explicitly defined.
   \item[line 19-28] defines some constants used in the main program
   \item[line 33] here the structure \texttt{meshgen\_options}
     of type
      \texttt{meshgen\_options\_t} is declared.
 This structure is used to fill in all kinds
 of options for mesh generation. Throughout \tfem\ the convention is used
 that all derived types have an extension \texttt{\_t} added. 
   \item[line 34 and 36] here the main structures are defined that are used
throughout \tfem: \texttt{mesh} and \texttt{problem}, which are of type
\texttt{mesh\_t} and \texttt{problem\_t}, respectively. The mesh structure
contains everything related to the mesh (topology, coordinates, \ldots) and the
problem structure contains everything related to the problem (number of
unknowns in the nodes, nodes with essential boundary conditions, \ldots).

\item[line 35] like \texttt{meshgen\_options}, \texttt{input\_probdef} is used
to fill in various data concerning the problem before entering the problem
definition routine that finally fills the structure \texttt{problem}.

\item[line 37] this line defines the system matrix \texttt{sysmatrix}. It is a
structure of type \texttt{sysmatrix\_t}.

\item[line 38] this line defines all system vectors: the solution vector, the
right-hand side and the exact solution.  They are all
structures of type \texttt{sysvector\_t}.

\item[line 39] here the structure \texttt{plot\_options} of type
\texttt{plot\_options\_t} is declared.
 This structure is used to fill in all kinds
 of options for plotting with figplot. 

\item[line 40] here the structure \texttt{coefficients} of type
\texttt{coefficients\_t} is declared.
 This structure is used to transfer parameters to the element routines, such
 as order of interpolation, material parameters etc.

\end{description}


\begin{listingcont}
   
! fill coefficients

  call create_coefficients ( coefficients, ncoefi=100, ncoefr=50 )

  coefficients%i = 0
  coefficients%i(1) = uintpl
  coefficients%i(10:12) = [ gauss, gaussb, funcnr ]

  coefficients%r(1) = alpha
  coefficients%r(2:) = 0

  coefficients%func => func

\end{listingcont}

Here the `coefficients' are filled for the Poisson problem. The coefficients
are defined using a structure of type \texttt{coefficients\_t} and are split
into an array for integer parameters \texttt{coefficients\%i} and real ('double
precision') parameters \texttt{coefficients\%r}. Furthermore, functions need to be
assigned to ``procedure pointers'' in structure \texttt{coefficients} in order to
make them available on element level, as
shown above for the user scalar function \texttt{func} on line 54..
See Sec.~\ref{sec:convdiff_elements} how the coefficients are defined for the
Poisson add-on.

\begin{listing}{54}

! create mesh
                                                                                
  meshgen_options%elshape = 6
  meshgen_options%nx = nx
  meshgen_options%ny = ny

  call quadrilateral2d ( mesh, meshgen_options )

  call fill_mesh_parts ( mesh ) 

\end{listing}
Here the mesh is created:
\begin{description}
   \item[line 61] this is the call to the routine 
\texttt{quadrilateral2d}, which is a simple mesh generator for quadrilateral
domains. The default is to use a unit square.
   \item[line 57] \texttt{elshape=6} defines a nine node quad element.
 The shape numbers of \sep\ are used as much as possible here.
   \item[line 58 and 59] \texttt{nx} and \texttt{ny} are the number of elements
in $x$ and $y$ direction respectively.
   \item[line 63] this routine computes other parts of \texttt{mesh} from the
     topology produced by the mesh generator. These parts include, for example,
     an array that gives the nodes that are connected to a certain node. That
means that if you produce the topology by an external mesh generator, you
still have to call this routine to fill the various parts of the mesh
structure.
\end{description}

\begin{listingcont}
! problem definition
                                                                                
  call create_input_probdef ( mesh, input_probdef )

  input_probdef%elementdof(1)%a = [1,1,1,1,1,1,1,1,1]

  call define_essential ( mesh, input_probdef, curve1=1, curve2=4 )

  call problem_definition ( input_probdef, mesh, problem )

\end{listingcont}
Here the structure \texttt{problem} is filled:
\begin{description}
   \item[line 73] on this line the structure \texttt{problem} is actually
filled from the structure \texttt{input\_probdef}. Before we can do that we have
to create and fill the structure \texttt{input\_probdef} with some data.
   \item[line 67] here the structure \texttt{input\_probdef} is created,
      which includes allocation of some memory. On line 67 the structure
     \texttt{input\_probdef} is defined, in a sense created, but it is still
     very empty.
   \item[line 69] here the degrees of freedom of the element are defined, in
     this case one unknown in each nodal point. This means we use a bi-quadratic
     interpolation ($Q_2$). Some explanation here: the
component \texttt{elementdof} is an array of size equal to the number of
element groups, which explains the \texttt{(1)}, since we have only one element
group. Furthermore, in \ftn\ it is not possible to define arrays of arrays. To
circumvent that problem we use arrays of a derived type, where the latter has a
single component: an array. We adopted the convention to call that component
\texttt{a}, which explains the \texttt{\%a}.
   \item[line 71] here we define where we have essential (Dirichlet)
boundary conditions. The mesh generator \texttt{quadrilateral2d} defines four
curves as depicted in Fig.~\ref{fig:unitsquare}. On all curves, that is from
\texttt{curve1=1} until \texttt{curve2=4} we define essential boundary
conditions. All degrees of freedom are made essential, by default, but there is
only one here in each node.
\end{description}

\begin{listingcont}
! create system vectors (solution and right-hand side)
                                                                                
  call create_sysvector ( problem, sol )
  call create_sysvector ( problem, rhsd )
                                                                                
! fill solution vector with essential boundary conditions
                                                                                
  call fill_sysvector ( mesh, problem, sol, &
    curve1=1, curve2=4, func=func, funcnr=3 )

\end{listingcont}
Here the system vectors are created and filled:
\begin{description}
   \item[line 77 and 78] the system vectors of type \texttt{sysvector\_t} are
created here. This means that by this routine the data is actually allocated.
It is however not initialized (it might contain rubbish). If the user wants a
vector to be initialized to zero, say \texttt{sol}, just add a statement
\texttt{sol\%u=0} after the creation of the vector. The component of the
structure \texttt{sol} containing the data is \texttt{u}.
Here we come to a basic
principle in the design of \tfem: each routine is as elementary as possible and
basically performs a single task which is well defined. This has the advantage
that the user has full control as much as possible and nothing happens behind
his/her back. A disadvantage is of course that main programs can be quite long.
   \item[line 82 and 83] here the solution vector \texttt{sol} is filled with
  the values generated by the function \texttt{func}
  specified using a function number of 3 
 along curves \texttt{C1} to \texttt{C4}. Note that \texttt{sol} needs to be
created first before the fill routine is called, otherwise an error message is
given. This conforms with our basic design principle discussed above.
\end{description}

\begin{listingcont}
! create system matrix
                                                                                
  call create_sysmatrix_structure ( sysmatrix, mesh, problem, symmetric=.true. )
                                                                                
  call create_sysmatrix_data ( sysmatrix )
                                                                                
\end{listingcont}
Here the system matrix of type \texttt{sysmatrix\_t} is created:
\begin{description}
   \item[line 87] here the structure of the system matrix is created, i.e.\ the
     number of non-zeros in each row of the sparse matrix. The matrix is
defined to be symmetric so that only the upper diagonal needs to be stored.
   \item[line 89] here the actual matrix elements are allocated. 
\end{description}

\begin{listingcont}
! build (assemble) matrix and vector from elements
                                                                                
  call build_system ( mesh, problem, sysmatrix, rhsd, elemsub=poisson_elem, &
    coefficients=coefficients )
                                                                                
  call add_effect_of_essential_to_rhs ( problem, sysmatrix, sol, rhsd )
                                                                                
  call solve_system_ma57 ( sysmatrix, rhsd, sol )
                                                                                
\end{listingcont}
Here the system matrix and vector is build and the system is solved:
\begin{description}
   \item[line 93] here the matrix and vector is assembled from
 the element
 subroutine \texttt{poisson\_elem} and the function \texttt{func} from the
module \texttt{poisson\_functions\_m}. 
Note, that
\texttt{build\_system} is really a dumb routine that just calls \texttt{elemsub}
and doesn't know anything what goes on in there. 
   In the element subroutine \texttt{poisson\_elem} of the
     module \texttt{poisson\_elements\_m} a 
     variable \texttt{funcnr} has been defined that is used as function number
     to the scalar function routine \texttt{func}. The actual value of
     \texttt{funcnr} will be set to the value in \texttt{coefficients\%i(12)}.
     The right-hand side $f$ will be set to outcome of the function.
     Here we see our first
     example of communicating with the element routines via coefficients.
   \item[line 96] the system of equations
\begin{equation}
   Su=f
\end{equation}
has been partitioned as follows
\begin{equation}
   \begin{pmatrix}
      S_{\text{uu}} & S_{\text{up}} \\
      S_{\text{pu}} & S_{\text{pp}} 
   \end{pmatrix} 
   \begin{pmatrix}
      u_\text{u} \\
      u_\text{p}
   \end{pmatrix}=
   \begin{pmatrix}
      f_\text{u} \\
      f_\text{p}
   \end{pmatrix}
   \label{eq:partitionedsystem}
\end{equation}
where subscript $_\text{u}$ and $_\text{p}$ means ``unknown'' and ``prescribed'',
respectively. The prescribed part represents the essential boundary conditions.
The components of the vectors of type \texttt{sysvector\_t} are renumbered and
stored accordingly to
represent this partitioning\footnote{Note that this is different from \sep\,
where the solution vectors retain there original mesh numbering storage.}. Also
the system matrix contains all four sub-matrices. From
Eq.~\eqref{eq:partitionedsystem} we find the equation to solve:
\begin{equation}
   S_{\text{uu}}u_\text{u}=f_{\text{u}}-S_{\text{up}}u_\text{p}
  \label{eq:systemtosolve}
\end{equation}
The subroutine \texttt{add\_effect\_of\_essential\_to\_rhs} computes the last
term and adds it to $f_{\text{u}}$.

\item[line 98] here the system Eq.~\eqref{eq:systemtosolve} is solved using the
sparse direct solver MA57 for a symmetric matrix 
 (see Sec.~\ref{sec:hslsparse}) \cite{HSL}. At the end of the routine the
$LU$-factors are deallocated from memory. It is possible to keep the factors
for a back-substitution only, the next time. We refer to the comments in the
source to learn how to do that. 
\end{description}

\begin{listingcont}
! create exact solution

  call create_sysvector ( problem, solexact )

  call fill_sysvector ( mesh, problem, solexact, &
    node1=1, node2=mesh%nnodes, func=func, funcnr=3 )

! print maximum difference of sol-solexact to standard output

  print *, maxval( abs(sol%u -solexact%u) )
                                                                                
\end{listingcont}
Here an exact solution in the nodal points is created and the maximum
absolute difference is printed between the numerical solution and the exact
one:
\begin{description}
   \item[line 104 and 105] \texttt{mesh\%nnodes} is the number of nodal points.
Note that the heading of \texttt{fill\_sysvector} has different parameters than
the call on line 82--83. This is possible due to a very nice feature of \ftn:
 optional parameters.
   \item[line 109] the array syntax of \ftn\ comes in handy here.
\end{description}

\begin{listingcont}
! plot solution

  call plot_color_contour ( plot_options, mesh, problem, filename='sol.fig', &
    sysvector=sol ) 

\end{listingcont}
This is an output section that prints plotting data to a file. Use
\texttt{.fig} as extension for plots made with figplot.

\begin{listingcont}
! delete all data including all allocated memory
                                                                                
  call delete ( problem )
  call delete ( input_probdef )
  call delete ( mesh )
  call delete ( sol, solexact, rhsd )
  call delete ( sysmatrix )
  call delete ( coefficients )

end program poisson1
\end{listingcont}
Here all allocated data is deleted again to end with a clean state.
In this case the deleting is just
clean programming and since the program stops after that, it is not absolutely
necessary. However it is recommended to always do this.
Usually programs are much more complicated
and it is important to keep track of the memory allocated to make the memory
footprint as small as possible. For example, if you would rewrite the program
to do a mesh refinement study by adding a loop, you must deallocate the memory
at the end of each cycle to avoid an `out of memory'.

The procedure \texttt{delete} is a so-called \emph{generic name} which
represents calls to other procedures (like \texttt{delete\_mesh},
\texttt{delete\_problem} etc.) depending on the type of the
arguments\footnote{There is a similar generic name \texttt{create}
for the routines beginning
with \texttt{create\_}, except for the routines that involve the system matrix.
However use this only if you are very experienced with \tfem, because the
various underlying \texttt{create\_} routine have a varying number and type of
arguments. Using only \texttt{create} it is less easy to jump to the source of
the routine using tags to check the required arguments.}
The number of arguments allowed has a maximum (at the time of writing five),
but all arguments must be of the same type.

Now compile/link the program \texttt{poisson1} by typing
\begin{quote}
   \texttt{make poisson1}
\end{quote}
or just
\begin{quote}
   \texttt{make}
\end{quote}
which compile/links all main programs in the current directory. Run the program
by
\begin{quote}
   \texttt{poisson1}
\end{quote}
and if all goes well you see a line on your screen similar to
\begin{quote}
\texttt{4.071378136849546E-07}
\end{quote}
The file \texttt{sol.fig} can be plotted by typing
\begin{quote}
   \texttt{viewfig sol.fig}
\end{quote}
which will show something like Fig.~\ref{fig:poisson1_sol.fig}.
\begin{figure}[ht!]
   \begin{center}
   \includegraphics[scale=0.5]{poisson1_sol}
   \end{center}
   \caption{Plot of \texttt{sol.fig} for \texttt{poisson1}}
    \label{fig:poisson1_sol.fig}
\end{figure}

\subsection{A Poisson problem on a unit square with Dirichlet
            and Neumann boundary conditions}

\label{sec:poisson2}
We consider the Poisson problem on a unit square, like in the previous section. 
We change the boundary condition on curve \texttt{C2} to a natural boundary,
meaning that we have to specify $h_N=-\plderiv{u}{n}$. 
We substitute the following exact solution for $u$:
\begin{equation}
     u=\cos(\pi x)\cos(\pi y) + x^3y^3
\end{equation}
and find the corresponding $f$:
\begin{equation}
     f=2\pi^2\cos(\pi x)\cos(\pi y) - 6(xy^3+x^3y)
\end{equation}
Specifying $-\plderiv{u}{n}$ is easy for the square domain, but for a
curvilinear domain it is very difficult. In that case, it is much easier to
specify $-\nabla u$ by a function
and let the boundary element compute $h_N=-\plderiv{u}{n}=-\vec n\cdot\nabla u$.
This is what we do here as well and specify $-\nabla u$ by a vector function
\texttt{vfunc}. The exact value for $\nabla u$ is given by
\begin{equation}
   \nabla u=
\begin{pmatrix}
     -\pi\sin(\pi x)\cos(\pi y) + 3x^2y^3 \\
     -\pi\cos(\pi x)\sin(\pi y) + 3x^3y^2 
\end{pmatrix}
\end{equation}
The source is in the file \texttt{poisson2.f90} and the difference with respect
to \texttt{poisson1.f90} will be discussed. The changes/added lines are:
\begin{listing}{22}
    funcnr=8,           & ! function number for the right-hand side
    vfuncnr=1             ! function number for the flux vector
\end{listing}
\bigskip
\begin{listing}{47}
  coefficients%i(6) = 2  ! use a vector function for h_N
\end{listing}
\bigskip
\begin{listing}{49}
  coefficients%i(16) = vfuncnr
\end{listing}
\bigskip
\begin{listing}{55}
  coefficients%vfunc => vfunc
\end{listing}
\bigskip
\begin{listing}{74}
  call define_essential ( mesh, input_probdef, curve1=1 )
  call define_essential ( mesh, input_probdef, curve1=3, curve2=4 )

\end{listing}
\bigskip
\begin{listing}{84}
! fill solution vector with essential boundary conditions
                                                                                
  call fill_sysvector ( mesh, problem, sol, &
    curve1=1, func=func, funcnr=7 )
  call fill_sysvector ( mesh, problem, sol, &
    curve1=3, curve2=4, func=func, funcnr=7 )

\end{listing}
\bigskip
\begin{listing}{102}
  call add_boundary_elements ( mesh, problem, rhsd, curve=2, &
    elemsub=poisson_natboun, coefficients=coefficients )

\end{listing}
\bigskip
\begin{listing}{113}
  call fill_sysvector ( mesh, problem, solexact, &
    node1=1, node2=mesh%nnodes, func=func, funcnr=7 )
                                                                                
\end{listing}
\begin{description}
   \item[line 22 and 23] changed \texttt{funcnr} and added \texttt{vfuncnr}.
   \item[line 47 and 49] set for vector function (line 47) and
     add \texttt{vfuncnr} to the coefficients. Note, that the
vector function needs to specify the flux vector $-\nabla u$.
   \item[line 55] the vector function \texttt{vfunc} is assigned to
      the procedure pointer \texttt{coefficients\%vfunc}.
   \item[line 74 and 75] we have split the essential boundaries into two parts:
\texttt{C1} and \texttt{C3--C4}.
   \item[lines 84--90] the filling has been split in the same way.
Note, that this
is not really necessary here, it is just clean programming.
Filling the values on
curve \texttt{C2} would have done no harm here, since the function is
defined on curve \texttt{C2} and the unknowns are recomputed by the
solve step (they are unknowns now!). In general this practice is not
recommended, since functions might not be defined.
   \item[line 102 and 103]here the natural boundary integral is added using the
element subroutine
\texttt{poisson\_natboun}
and using the vector function \texttt{vfunc}. The function number
is supplied via the coefficients.
    \item[line 113 and 114]
      the exact solution must now be build with \texttt{funcnr=7}.
\end{description}

Typing at the command line 
\begin{quote}
    \texttt{poisson2}
\end{quote}
will show something like
\begin{quote}
  \texttt{4.444285222060529E-07}
\end{quote}
on your screen and typing
\begin{quote}
  \texttt{viewfig sol.fig}
\end{quote}
will show something like Fig.~\ref{fig:poisson2_sol.out}.
\begin{figure}[ht!]
   \begin{center}
   \includegraphics[scale=0.5]{poisson2_sol}
   \end{center}
   \caption{Plot of data in \texttt{sol.fig} for \texttt{poisson2}}
    \label{fig:poisson2_sol.out}
\end{figure}

\subsection{A Poisson problem on a curved quadrilateral region with Dirichlet 
            and Neumann boundary conditions} 
\label{sec:poisson3}

We again consider the Poisson equation as in the previous example, but now we
consider a quadrilateral domain with curved boundaries.
We substitute the following exact solution for $u$:
\begin{equation}
     u=1+\cos(\pi x)\cos(\pi y)
\end{equation}
and find the corresponding $f$:
\begin{equation}
     f=2\pi^2\cos(\pi x)\cos(\pi y)
\end{equation}
and $\nabla u$:
\begin{equation}
   \nabla u=
\begin{pmatrix}
     -\pi\sin(\pi x)\cos(\pi y) \\
     -\pi\cos(\pi x)\sin(\pi y) 
\end{pmatrix}
\end{equation}
The source of the problem is in the file \texttt{poisson3.f90} and the
difference with respect to \texttt{poisson2.f90} will be discussed. 
The function numbers need to be changed again and the main change in the
source is
\begin{listing}{59}
  meshgen_options%regionshape = 3
  meshgen_options%x2d = &
    reshape ( [ 0.0_dp, 2.0_dp, 2.0_dp, -1.0_dp,    &
                0.0_dp, 0.0_dp, 2.0_dp,  3.0_dp ], &
              [4,2] )
  meshgen_options%curved(1:4) = .true.
  meshgen_options%funcnr(1:4) = [ 1, 2, 3, 4 ]
                                                                                
  call quadrilateral2d ( mesh, meshgen_options, func=cfunc )
   
\end{listing}
\begin{description}
   \item[line 59] \texttt{regionshape=3} is a curved quadrilateral
     (1: rectangle, 
     2: quadrilateral with straight boundaries)
   \item[line 60--63] here the four corner points are specified: $(0,0)$,
         $(2,0)$, $(2,2)$, $(-1,3)$. Note that \texttt{x2d} is a 4$\times$2
        matrix.
   \item[line 64] here it is specified that all four curves are curved. It is
      of course also possible to define only one, two or three curves curved.
   \item[line 65] here the function numbers for the $(x(s),y(s))$
       coordinates of the four curved boundaries as a function of a curve
       coordinate $s\in[0,1]$. The curve coordinate $s$ increases in the
       positive direction of a curve, as given in Fig.~\ref{fig:unitsquare}.
       The start point ($s=0$) and end point
       ($s=1$) of a curve 
       need to coincide with the corner coordinates \texttt{x2d}, otherwise
       weird results can be expected.
   \item[line 67] an extra parameter specifying the function for the curved
       boundaries is present.
       \texttt{cfunc} is part of the module \texttt{functions\_m}.
\end{description}

Typing at the command line 
\begin{quote}
    \texttt{poisson3}
\end{quote}
will show something like
\begin{quote}
  \texttt{8.482015264414944E-04}
\end{quote}
on your screen and typing
\begin{quote}
  \texttt{viewfig}
\end{quote}
will show something like
Figs.~\ref{fig:poisson3_mesh.fig} and \ref{fig:poisson3_sol.fig}. 
\begin{figure}[ht!]
   \begin{center}
     \includegraphics[scale=0.7]{poisson3_mesh}
   \end{center}
   \caption{Plot of the mesh for \texttt{poisson3} (20$\times$20 elements).}
    \label{fig:poisson3_mesh.fig}
\end{figure}
\begin{figure}[ht!]
   \begin{center}
   \includegraphics[scale=0.7]{poisson3_sol}
   \end{center}
   \caption{Plot \texttt{sol.fig} for \texttt{poisson3}}
    \label{fig:poisson3_sol.fig}
\end{figure}
Note, that all elements are curved and not just the boundary. The mesh is
computed by a mapping technique called `the blending function method' of Gordon
\& Hall. This technique is fully described in the book of
 Szab{\'o} and Babu{\v{s}}ka \cite{Szabo91}. First a mesh is made on the unit
square $(\xi,\eta)\in[0,1]\times[0,1]$. Then the mesh is mapped on a
quadrilateral
with straight sides (bilinear function). Now assume that, for example, curve
\texttt{C2} is curved. Now all points are shifted in the following way. Say we
have a particular point denoted by the reference coordinate $(\xi_1,\eta_1)$.
Now take the point on the boundary \texttt{C2}: $(1,\eta_1)$. The new position
of that point is computed using the boundary description with $s=\eta_1$. Then
the difference with the position of the boundary point in the straight
quadrilateral is computed.
The new position of the internal point given by reference coordinates
$(\xi_1,\eta_1)$
is found by adding the
same difference but weighted (blended) by $\xi$. Note, that the other sides
remain straight. The same operation can be performed for all sides that are
curved.

\subsection{A Poisson problem on a curved quadrilateral region with Dirichlet 
            and Neumann boundary conditions. Derived quantities.} 
\label{sec:poisson4}

We again consider the Poisson equation as in the previous example, using the
same exact solution. However, we want to compute the $\nabla u$ from the
numerical solution $u$. This is a derived quantity. 

For derived quantities (but also for other purposes) we need a new vector type.
The type \texttt{sysvector\_t} is exclusively for the system vector:
\begin{itemize}
   \item it is governed by the array component \texttt{elementdof} of
      the structure \texttt{input\_probdef\%elementdof}.
   \item it is also renumbered to reflect unknown and known entries in the
      vector.
\end{itemize}
Derived quantities usually have a different number of components in the nodal
points than the system vector and there is no reason to renumber the
components. Therefore we introduce the new vector type \texttt{vector\_t} and
a number of associated routines. Vectors of type \texttt{vector\_t} are
structures like system vectors and can be defined like a system vector without
the renumbering using \texttt{vec\_elementdof} (see below in the example). An
additional possibility exists to define the components element-wise. That is,
the degrees defined by \texttt{vec\_elementdof} are defined per element so that
each element connected to a particular nodal point has its own set of
data for the
vector components at that point. This can come in handy, for example, when the
user wants to see the `real' derivatives of a \fem\ solution and not just
a nodal averaged one.

The source of the problem is in the file \texttt{poisson4.f90} and the
difference with respect to \texttt{poisson3.f90} will be discussed. 
\bigskip
\begin{listing}{35}
  type(sysvector_t), target :: sol
  type(sysvector_t) :: rhsd
  type(vector_t) :: grad, gradexact
  type(plot_options_t) :: plot_options
  type(oldvectors_t) :: oldvectors
  type(coefficients_t) :: coefficients
\end{listing}
\bigskip
\begin{listing}{74}
! problem definition
                                                                                
  call create_input_probdef ( mesh, input_probdef, nvec=1 )
                                                                                
  input_probdef%elementdof(1)%a = [1,1,1,1,1,1,1,1,1]

  input_probdef%vec_elementdof(1)%a = reshape ( [2,2,2,2,2,2,2,2,2], [9,1])
\end{listing}
\bigskip
\begin{listing}{117}
! create of grad
                                                                                
  call create_vector ( problem, grad, vec=1 )

! create oldvectors

  call create ( oldvectors, nsysvec=1 )
                                                                                
  oldvectors%s(1)%p => sol

  call derive_vector ( mesh, problem, grad, elemsub=poisson_deriv, &
    coefficients=coefficients, oldvectors=oldvectors )
                                                                                
! create exact solution of grad
                                                                                
  call create_vector ( problem, gradexact, vec=1 )
                                                                                
  call fill_vector ( mesh, problem, gradexact, degfd=1, &
    node1=1, node2=mesh%nnodes, func=func, funcnr=5 )
  call fill_vector ( mesh, problem, gradexact, degfd=2, &
    node1=1, node2=mesh%nnodes, func=func, funcnr=6 )
                                                                                
! print maximum difference of grad-gradexact to standard output
                                                                                
  print *, maxval( abs(grad%u -gradexact%u) )
                                                                                
! plot solution

  plot_options%fontsize = 14

  call plot_color_contour ( plot_options, mesh, problem, filename='dudx.fig', &
    degfd=1, vector=grad )
  call plot_color_contour ( plot_options, mesh, problem, filename='dudy.fig', &
    degfd=2, vector=grad )

  plot_options%scalevector = 0.05

  call plot_vector ( plot_options, mesh, problem, filename='gradu.fig', &
    vector=grad )

\end{listing}
\bigskip
\begin{listing}{163}
  call delete ( grad, gradexact )
  call delete ( sysmatrix )
  call delete ( coefficients )
  call delete ( oldvectors )
\end{listing}

\begin{description}
   \item[line 35] the solution vector \texttt{sol} gets the additional
attribute \texttt{target} so that a `pointer' in
\texttt{oldvectors} can point to it (see below).
   \item[line 37] defines two vectors of type \texttt{vector\_t}.
   \item[line 39] defines the structure \texttt{oldvectors} of
      type \texttt{oldvectors\_t}. This structure is used to supply sysvectors
      and vectors to the element routines by using pointers.
   \item[line 76] additional parameter \texttt{nvec=1} to indicate that there
     is one subtype of vector defined.
     Here is not meant \texttt{vector\_t} but
     the number of vectors with a different number of degrees of freedom in the
     nodal points. For
     example, a vector with one degree of freedom in all nodal points and a
     vector with two degrees of freedom in all nodal points are both of type
     \texttt{vector\_t} but of different subtype. 
  \item[line 80] here all subtypes of the vectors are defined. In this case
     there is only one: a vector with two degrees of freedom in each nodal
     point. Note, that \texttt{vec\_elementdof} is an array of size the number
     of element groups. Since we have only one group, only one array
     component (hence the \texttt{(1)}) needs to be defined. Each component of
     the array is a \emph{two-dimensional} array, which is denoted by
     \texttt{\%a}. The first dimension is the number of nodal points in the
     element and the second dimension is the vector subtype number. We have
     only one subtype here, but it still needs to be defined as a
    two-dimensional array, for which we have used the intrinsic function
      \texttt{reshape}.
   \item[line 119] here a vector \texttt{grad}
     of type \texttt{vector\_t} with subtype number 
     \texttt{vec=1} is created. The subtype number is stored in the component
     \texttt{grad\%vec}.
   \item[line 123] here the structure oldvectors is defined (it contains just
     one sysvector).
   \item[line 125] here the solution vector \texttt{sol} is communicated to the
     element routine \texttt{poisson\_deriv} by using a pointer. The
     assignment statement here does not transfer the data of \texttt{sol} to
     \texttt{oldvectors\%s(1)\%p}:
     \emph{only the pointer} \texttt{oldvectors\%s(1)\%p}
     \emph{will be associated to} the system vector \texttt{sol}.
     This means that after this line \texttt{oldvectors\%s(1)\%p} points to
     the same object (including the data) in memory as \texttt{sol}.
   \item[line 127] the actual values in \texttt{grad} are computed
     here using the element subroutine \texttt{poisson\_deriv}. The numerical
     values of $\nabla u$ will be discontinuous across element boundaries.
     Therefore the values on element boundaries will be averaged values.
     It is possible to define the vector \texttt{grad} to be defined
     element-wise, in which case averaging will not take place.
   \item[line 132--137] here the vector with the exact $\nabla u$ is
     created and filled.
   \item[line 141--155] the maximum error in $\nabla u$ is printed to
     the screen and \texttt{gradu} is plotted to a file.
   \item[line 163] the vectors are deleted to free the memory.
   \item[line 166] the structure oldvectors is deleted.
\end{description}
     
Typing at the command line 
\begin{quote}
    \texttt{poisson4}
\end{quote}
will show something like
\begin{quote}
  \texttt{9.612474932170202E-02}
\end{quote}
on your screen and typing
\begin{quote}
  \texttt{viewfig}
\end{quote}
will show something like
Figs.~\ref{fig:poisson4_1_sol.fig}, \ref{fig:poisson4_2_sol.fig}, 
      and \ref{fig:poisson4_3_sol.fig}.
Note that, compared to the previous problem, we have made curve \texttt{C2}
a straight line.
\begin{figure}[ht!]
   \begin{center}
   \includegraphics[scale=0.7]{poisson4_1_sol}
   \end{center}
   \caption{Plot of $\plderiv{u}{x}$ for \texttt{poisson4}.}
    \label{fig:poisson4_1_sol.fig}
\end{figure}
\begin{figure}[ht!]
   \begin{center}
   \includegraphics[scale=0.7]{poisson4_2_sol}
   \end{center}
   \caption{Plot of $\plderiv{u}{y}$ for \texttt{poisson4}.}
    \label{fig:poisson4_2_sol.fig}
\end{figure}
\begin{figure}[ht!]
   \begin{center}
   \includegraphics[scale=0.7]{poisson4_3_sol}
   \end{center}
   \caption{Vector plot of $\nabla u$ for \texttt{poisson4}.}
    \label{fig:poisson4_3_sol.fig}
\end{figure}

\clearpage

\section{The Stokes equation}

In this section an example using the Stokes equation will be given after
the theory in the next section.

\subsection{Weak form of the Stokes equation. Galerkin \fem}
\label{sec:stokesweak}

We consider the Stokes equation in a region $\Omega$:
\begin{alignat}{3}
 -&\nabla\cdot(2\eta \ten D)+\nabla p & &=\vek f,&&\qquad\text{in }\Omega
     \label{eq:stokes1}\\
   &\nabla\cdot\vek u& &=0&&\qquad\text{in }\Omega
     \label{eq:stokes2}
\end{alignat}
where $\ten D=(\nabla\vek u+\nabla\vek u^T)/2$. The boundary conditions are
\begin{gather}
    \vec u = \vec u_D\text{ on }\Gamma_D\\
    -p\vec n+2\eta \ten D\cdot\vec n=\vec t_{\text{N}}
\text{ on }\Gamma_{\text{N}}
\end{gather}
where the boundary $\Gamma=\Gamma_D\cup\Gamma_{\text{N}}$ has been split into a
Dirichlet and a Neumann part. The Neumann condition is equivalent
to an imposed traction of $\vec t_{\text{N}}$.
                                                                                
The weak form can be obtained by multiplying
Eqs.~\eqref{eq:stokes1}--\eqref{eq:stokes2}
with test functions $\vec v$ and $q$:
\begin{align*}
 (\vek v,-\nabla\cdot\ten\sigma -\vek f)&=0,\qquad\text{for all }\vek v\\
  (q,\nabla\cdot\vek u)& =0,\qquad\text{for all }q
\end{align*}
where
\begin{align*}
   (\vek a, \vek b)&=\int_\Omega\vek a\cdot\vek b\D x\\
   (a,b)&=\int_\Omega a b\D x
\end{align*}
and $\ten\sigma = -p\ten I+2\eta\ten D$.

Using partial integration and Gauss'
theorem we obtain the following weak form: 
Find $\vek u$ and $p$ such that
\begin{alignat}{3}
  &\big(\ten D_v,2\eta\ten D\big)
    -(\nabla\cdot\vek v,p)&&=(\vek v,\vek t_{\text{N}})_{\Gamma_{\text{N}}}+(\vek v,\vek f)
&\qquad&\text{for all }\vek v\label{eq:weakst1}\\
  &\,(q,\nabla\cdot\vek u)&& =0,&\qquad&\text{for all }q\label{eq:weakst2}
\end{alignat}
using appropriate spaces for $\vek u$, $p$, $\vek v$ and $q$. In this equation
$\ten D_v$ is defined by
\[
\ten D_v=(\nabla\vek v+\nabla\vek v^T)/2
\]
and the bilinear operator $(\,,\,)$ for tensors is given by
\[
   (\ten A, \ten B)=\int_\Omega\ten A:\vek B\D x
\]

The weak form can be used to obtain an approximate solution using the Galerkin \fem\
technique. For this we divide the region into elements:
\[
\Omega=\cup_e\Omega_e, \quad\Omega_{e}\cap\Omega_{f}=\emptyset\quad
   \text{for }e\neq f
\]
and interpolate $\vek u$ and $p$ (and similarly $\vek v$ and $q$) by
polynomials on each element:
\[
 \begin{split}
    \vek u_h&=\tsum_k \phi_k(\vek x)\vek u_k=\col\phi^T(\vek x)\col{\vek u}\\
    p_h&=\tsum_k \psi_k(\vek x) p_k=\col\psi^T(\vek x)\col{ p}
 \end{split}
\]
where $\col\phi(\vek x)$ and $\col\psi(\vek x)$
are global shape functions.  Substituting into the
weak form leads to
\begin{alignat*}{3}
   &\col{\vek v}^T\cdot(\mat{\ten S}\cdot\col{\vek u}-\mat{\vek L}^T\col p)&&=
    \col{\vek v}^T\cdot\col{\vek f}\qquad&&\text{for all }\col{\vek v}\\
  &\quad\col q^T\mat{\vek L}\cdot\col{\vek u}&&=0\qquad&&\text{for all }\col{q}
\end{alignat*}
or
\[
 \begin{bmatrix}
   \mat{\ten S} &-\mat{\vek L}^T \\[1ex]
   \mat{\vek L} &\quad \mat 0
 \end{bmatrix}
 \begin{bmatrix}
    \col{\vek u}\\[1ex]
    \col p
 \end{bmatrix}=
 \begin{bmatrix}
    \col{\vek f}\\[1ex]
    \col 0
 \end{bmatrix}
\]
with 
\begin{align*}
   \mat{\ten S}&=\int_\Omega\eta[(\nabla\col{\phi}\cdot\nabla\col{\phi}^T)\ten
I+                            (\nabla\col{\phi}\nabla\col{\phi}^T)^c]\D x,
\qquad\text{where }\ten A^c\equiv\ten A^T\\
\mat{\vek L}&=(\col{\psi},\nabla\col\phi^T)\\
 \col{\vek f}&=(\col{\phi},\vek t_{N})_{\Gamma_{\text{N}}}+(\col{\phi},\vek f)
\end{align*}
or in index notation:
\begin{align*}
       S_{ij}^{km}&= \tsum_l(\phi_{k,l},\eta\phi_{m,l})\delta_{ij}+
          (\phi_{m,i},\eta\phi_{k,j})\\
    L_i^{km}&=(\psi_k,\phi_{m,i})\\
     f^k_i&=(\phi_k,t_{iN})_{\Gamma_{\text{N}}}+(\phi_k,f_i)
\end{align*}
Here $k$ and $m$ denote nodal points and $i$ and $j$ component directions of
vectors. In practice we multiply the continuity equation by -1 in order to
obtain a symmetric system matrix:
\[
 \begin{bmatrix}
   \mat{\ten S} &-\mat{\vek L}^T \\[1ex]
   -\mat{\vek L} &\mat 0
 \end{bmatrix}
 \begin{bmatrix}
    \col{\vek u}\\[1ex]
    \col p
 \end{bmatrix}=
 \begin{bmatrix}
    \col{\vek f}\\[1ex]
    \col 0
 \end{bmatrix}
\]

Recommended shape functions for velocity and pressure are:
\begin{itemize}
 \item from the Taylor-Hood
family (continuous pressures): $Q_2Q_1$ on quadrilaterals and hexahedra and
$P_2P_1$ on triangles and tetrahedra\footnote{The shape functions $P_2Q_2/P_1Q_1$ on prisms seem to also work in early testing (see example \texttt{stokes51.f90} in Section~\ref{sec:stokes3Ddrivencavity}).
However, the shape functions $P_2Q_2\text{r}/P_1Q_1\text{r}$ on pyramids do not seem to work. More research is needed here.}.
 \item from the Crouzeix-Raviart
family (discontinuous pressures): $Q_2P_1^{(\text{d})}$ on quadrilaterals
and hexahedra and $P_2^+P_1^{(\text{d})}$ on triangles and tetrahedra.
\item spectral elements (discontinuous pressures): $Q_pP_{p-1}^{(\text{d})}$ on
quadrilaterals.
\end{itemize}

For plotting streamlines in 2D flows we compute the streamfunction 
$\psi$ from the
velocity vector $\vec u=u\vec e_1+v\vec e_2$. The definition of $\psi$ is
\[
   \pderiv{\psi}{x} = -v,\qquad
   \pderiv{\psi}{y} = u
\]
which shows that
\begin{equation}
  \psi(\vek x_2) - \psi(\vek x_1) = \int_{\Gamma_{12}}
      \vek n\cdot\vek u\D s
    \label{eq:streamint}
\end{equation}
or in words: the streamfunction difference is the flow rate through a curve
connecting two points. Eq.~\eqref{eq:streamint} can be used to used to compute 
$\psi$ in all nodal points, but the procedure is cumbersome and not
straightforward. Therefore we prefer a procedure that solves 
the Poisson equation
\begin{equation}
   -\nabla^2\psi = \omega
\label{eq:streamfunction} 
\end{equation}
with the vorticity $\omega=\plderiv vx-\plderiv uy$. We use a
Neumann boundary condition:
\begin{equation}
 \begin{split}
   \pderiv{\psi}{n}=\vec n\cdot\nabla\psi=\vec n\cdot(-v\vec e_1+u\vec e_2)=
        -vn_1+un_2=-\vec u\cdot\vec t
 \end{split}
\label{eq:streamNeum}
\end{equation}
where $\vec t$ is the tangential vector (with $\vec n$ to the `right').

For plotting streamlines in an axisymmetrical flow we have
\begin{equation}
   \pderiv{\psi}{z} = -2\pi r v,\qquad
   \pderiv{\psi}{r} = 2\pi r u
   \label{eq:streamfuncaxi}
\end{equation}
where $(z,r)$ are cylindrical coordinates\footnote{We use the convention
$(z,r)$, i.e.\ $z$ is the first coordinate, to preserve the ordering of the
coordinates to represent a right-handed system.
Note that \sep\ uses the convention $(r,z)$.}
and $u$ is
the velocity component in axial ($z$) direction and $v$ the velocity component
in radial ($r$) direction ($\vec u=u\vec e_z+v\vec e_r$).
We easily find that
\[
  \psi(\vek x_2) - \psi(\vek x_1) = \int_{\Gamma_{12}}
      \vek n\cdot\vek v\,2\pi r\D s
\]
or in words: the streamfunction difference is the flow rate through a surface
consisting of a curve 
connecting two points that is expanded over the full circumferential ($\theta$)
direction.
We find from \eqref{eq:streamfuncaxi} that
\[
   -\nabla^2_{zr}\psi=-(\pderiv{^2\psi}{z^2}+\pderiv{^2\psi}{r^2})=
    = 2\pi ( r\omega+u)
\]
with the vorticity $\omega=\plderiv vz-\plderiv ur$. Note that
$\nabla^2_{zr}$ is the `planar' $(z,r)$ Laplace operator and not the 
real $\nabla^2$ expressed in a cylindrical coordinate system.
We use a Neumann boundary condition:
\begin{equation}
 \begin{split}
   \pderiv{\psi}{n}=\vec n\cdot\nabla\psi=
       2\pi r\vec n\cdot(-v\vec e_z+u\vec e_r)=
       2\pi r( -vn_z+un_r )=-2\pi r\vec u\cdot\vec t
 \end{split}
\label{eq:streamNeumaxi}
\end{equation}
where $\vec t$ is the tangential vector (with $\vec n$ to the `right'). Note,
that the Neumann condition is zero on the center line ($r=0$).

\subsection{A Stokes problem on a unit square with Dirichlet boundary 
            conditions: driven cavity problem. Physical quantities}
\label{sec:driven_cavity2D}

We consider the Stokes equation for a driven cavity problem on a unit square
(see Fig.~\ref{fig:drivencavity}). On the upper boundary a tangential
velocity of 1 is imposed and on all other walls the velocity components are
zero. In the lower left corner the pressure value is set to zero (Dirichlet).
\begin{figure}[ht!]
   \begin{center}
   \includegraphics{driven_cavity}
   \end{center}
   \caption{Driven cavity problem}
    \label{fig:drivencavity}
\end{figure}

The example file is \texttt{stokes1.f90} and can be obtained by typing at the
command line:
\begin{quote}
   \texttt{getexample stokes}
\end{quote}

The Stokes element routines
in the add-on `stokes' (see Sec.~\ref{sec:stokesaddon}) are used for
the building of the system. 

Some notes on the source:
\begin{listing}{43}

! fill coefficients

  call create_coefficients ( coefficients, ncoefi=150, ncoefr=100 )

  coefficients%i(1:11) = &
    [ uintpl,   pintpl,     0,     0,         0,  &
      physqvel, physqpress, 0,     0,     gauss,  &
      gauss ]
  coefficients%i(12:) = 0

  coefficients%r(1) = eta
  coefficients%r(2:) = 0

\end{listing}
Here the structure \texttt{coefficients} of type \texttt{coefficients\_t} (see
Sec.~\ref{sec:coefficients}) is created and filled with various parameters
needed on element level for the Stokes problem. See Sec.~\ref{sec:stokesaddon}
on how obtain the definition of the values in \texttt{coefficients} for the
Stokes system.
\begin{listing}{76}
! problem definition
                                                                                
  call create_input_probdef ( mesh, input_probdef, nvec=3, nphysq=2 )
                                                                                
  input_probdef%vec_elementdof(1)%a =   &
      reshape ( [2,2,2,2,2,2,2,2,2,    &  ! velocity
                 1,0,1,0,1,0,1,0,0,    &  ! pressure
                 1,1,1,1,1,1,1,1,1 ], &  ! scalar, such as vorticity
                 [9,3] )
                                                                                
  input_probdef%physq = [1,2]
                                                                                
  call define_essential ( mesh, input_probdef, curve1=1, curve2=4, physq=1 )
  call define_essential ( mesh, input_probdef, point=1, physq=2 )
                                                                                
  call problem_definition ( input_probdef, mesh, problem )
                                                                                
\end{listing}
\begin{description}
   \item[line 80--84] here three subtypes of vectors are defined. The first two
will also be used to define the degrees of freedom of the system vector by
using physical quantities (see below). We use bi-quadratic ($Q_2$)
and bilinear interpolation ($Q_1$) for the velocity vector and pressure,
respectively (Taylor-Hood element).
   \item[line 78] an extra parameter \texttt{nphysq=2} is added, which means
that we define two physical quantities (velocity and pressure).
   \item[line 86] here we define the two physical quantities, The first
(velocity) is defined by the vector of subtype \texttt{vec=1} and the second by
\texttt{vec=2}. We have not defined \texttt{input\_probdef\%elementdof}. This
means
that the element degrees of freedom for the system vector will automatically be
defined by the (sum of the) physical vectors.
   \item[line 88 and 89] essential boundary conditions can also be defined by
specifying the physical quantity instead of the degree of freedom number in the
nodes. Note that the pressure boundary condition is defined specifying
 the point number.
\end{description}
\begin{listing}{98}
! fill solution vector with essential boundary conditions
                                                                                
  call fill_sysvector ( mesh, problem, sol, &
    curves=[1,2,4], physq=1, value=0._dp )
  call fill_sysvector ( mesh, problem, sol, &
    curve1=3, physq=1, vvalue=[1._dp,0._dp] )
  call fill_sysvector ( mesh, problem, sol, &
    point=1, physq=2, value=0._dp )
                                                                                
\end{listing}
Here the actual values of the essential boundary conditions are filled in. Note
that in the upper corner nodes the velocity component $u$ has been set to 1. If
you want to exclude these end points add the parameter \texttt{exclude=3} at
the filling of curve \texttt{C3}.
We set $p=0$ in the lower left corner point.

\begin{listing}{113}
! build (assemble) matrix and vector from elements
                                                                                
  call build_system ( mesh, problem, sysmatrix, rhsd, elemsub=stokes_elem, &
    order = 'DN' )
                                                                                
\end{listing}
Here a new parameter \texttt{order} has been introduced. In the elements for
the Stokes problem the element degrees of freedom are not ordered node for node
(like in the global vector) but first all $u$ components, then all $v$
components and finally all $p$ components. This is reflected in the parameter 
\texttt{order='DN'}, which means ordering according to physical quantities
(outer loop) then degrees within physical quantities and finally
nodes (inner loop). In \texttt{revision 128} of \tfem\ the default has been
changed to \texttt{order='DN'} for problems with physical unknowns and the
parameter \texttt{order='DN'} is superfluous here.

\begin{listing}{122}
! post-processing
                                                                                
  call create_vector ( problem, velocity, physq=1 )
  call create_vector ( problem, pressure, vec=3 )
  call create_vector ( problem, vorticity, vec=3 )
                                                                                
  call extract_physvector ( mesh, problem, sol, velocity )

! create the structure oldvectors

  call create_oldvectors ( oldvectors, nsysvec=1 )

  oldvectors%s(1)%p => sol

  call derive_vector ( mesh, problem, pressure, elemsub=stokes_pressure, &
    coefficients=coefficients, oldvectors=oldvectors )

  coefficients%i(13)=5

  call derive_vector ( mesh, problem, vorticity, elemsub=stokes_deriv, &
    coefficients=coefficients, oldvectors=oldvectors )
                                                                                
\end{listing}
\begin{description}
   \item[line 124] a vector for the velocity is defined using
     \texttt{physq} instead of \texttt{vec} (both ways are equivalent).
   \item[line 125] a vector for the pressure is defined using \texttt{vec=3}.
Note that this vector defines the pressure in all nodes, whereas the pressure
in the solution vector \texttt{sol} is only defined in the vertices of the
element. The reason for this is that the plotting routines can only handle
vectors that have a value in all nodes.
   \item[line 128] the physical vector is extracted directly from the
     system vector.
   \item[line 132] here a structure of type \texttt{oldvectors\_t} is created
(see Sec.~\ref{sec:oldvectors}) to store the pointers to the solution that can
be used on element level. Note, that only one system vector is created in
\texttt{oldvectors}.
   \item[line 134] the pointer \texttt{oldvectors\%s(1)\%p} now points to
the solution vector \texttt{sol}.
   \item[line 136, 137 and 141,142] the pressure and vorticity vectors are
     computed here from the solution vector \texttt{sol}.
    The element routine for the
pressure \texttt{stokes\_pressure} computes the pressure in nine nodal points
from the four vertex values using the shape functions for the pressure.
\end{description}
Run the problem by typing
\begin{quote}
   \texttt{stokes1}\\
   \texttt{streamfunction1}
\end{quote}
The program \texttt{stokes1} also writes the velocity vector to a file, which is
read by the program \texttt{streamfunction1}. The streamfunction is computed by
solving the Poisson problem Eq~\eqref{eq:streamfunction} with the Neumann
boundary conditions Eq.~\eqref{eq:streamNeum}. Typing the command
\begin{quote}
  \texttt{viewfig}
\end{quote}
will show something like
Figs.~\ref{fig:streamfunction}, \ref{fig:velocity}, \ref{fig:vorticity}, 
      and \ref{fig:pressure} (and more).
Note, that due to the non-zero velocity in the
upper corner nodes we have some inflow and outflow near these nodes in the last
element. This can
be avoided by setting a zero velocity in these nodes. Either way, a singularity
exist in both points, best seen in the vorticity and pressure plots. 
\begin{figure}[ht!]
   \begin{center}
   \includegraphics[scale=0.5]{streamfunction}
   \end{center}
   \caption{Plot of the streamfunction for the driven cavity problem}
    \label{fig:streamfunction}
\end{figure}
\begin{figure}[ht!]
   \begin{center}
   \includegraphics[scale=0.5]{velocity}
   \end{center}
   \caption{Plot of the velocity vector for the driven cavity problem}
    \label{fig:velocity}
\end{figure}
\begin{figure}[ht!]
   \begin{center}
   \includegraphics[scale=0.5]{vorticity}
   \end{center}
   \caption{Plot of the vorticity for the driven cavity problem}
    \label{fig:vorticity}
\end{figure}
\begin{figure}[ht!]
   \begin{center}
   \includegraphics[scale=0.5]{pressure}
   \end{center}
   \caption{Plot of the pressure for the driven cavity problem}
    \label{fig:pressure}
\end{figure}
\clearpage

